\chapter{1911:Marie Curie}
\label{chap:chemistry1911}

The Nobel Prize in Chemistry for the year 1911 was awarded to Marie Curie.

\textbf{Motivation:}
in recognition of her services to the advancement of chemistry by the discovery of the elements radium and polonium, by the isolation of radium and the study of the nature and compounds of this remarkable element

\section{Marie Curie}

\subsection*{Early Life and Education}

Marie Curie was born on 1867-11-07 in Warsaw, Russian Empire (now Poland).
Maria Salomea Skłodowska Curie, better known as Marie Curie was a Polish and naturalised-French physicist and chemist. She shared the 1903 Nobel Prize in Physics with her husband, Pierre Curie, "for their joint researches on the radioactivity phenomena discovered by Professor Henri Becquerel". She won the 1911 Nobel Prize in Chemistry "for the discovery of the elements radium and polonium, by the isolation of radium and the study of the nature and compounds of this remarkable element".

\subsection*{Career and Major Research}

Curie worked at the Sorbonne in Paris, becoming the first woman to hold a professorship there. In 1898 she and her husband Pierre Curie discovered polonium and radium, and she later isolated pure radium metal, a feat of painstaking chemical work. She shared the 1903 Nobel Prize in Physics with Pierre Curie and Henri Becquerel for their studies of radioactivity.

\subsection*{Nobel Prize-winning Work}

The work for which Marie Curie was awarded the Nobel Prize in Chemistry in 1911 was described by the Nobel Committee as: "in recognition of her services to the advancement of chemistry by the discovery of the elements radium and polonium, by the isolation of radium and the study of the nature and compounds of this remarkable element".

This research fundamentally advanced the field by making groundbreaking contributions that transformed chemical science.

\subsection*{Significance and Impact}

The impact of Marie Curie's work extends far beyond the specific achievement recognized by the Nobel Prize.
These contributions continue to influence chemical research and education today.

\subsection*{Later Career and Honors}

Marie Curie continued her scientific work until 1934-07-04, passing away in Sallanches, France.
She received numerous honors including membership in major scientific academies and societies.

\subsection*{Legacy}

The legacy of Marie Curie endures through the lasting impact of her scientific contributions.
Marie Curie's work continues to be cited and built upon by researchers across the chemical sciences.

\subsection*{Modern Developments}

Marie Curie's discovery of radium and polonium established the science of radioactivity and gave medicine the tools of radiotherapy and radioactive tracers. Today, radiopharmaceuticals and techniques such as PET imaging and brachytherapy are direct descendants of her work. Her methods for isolating radioactive elements also laid the groundwork for the modern nuclear fuel cycle and radiation standards.

\subsection*{Selected Publications}

Marie Curie's major publications include numerous papers in leading journals of the era, detailing the experimental and theoretical work summarized above.

\clearpage

\section*{Chapter Summary}

Marie Curie was honored for the discovery of the elements radium and polonium, the isolation of radium, and the study of the nature and compounds of this remarkable element.
